% Folder 77, batch 07 — archivists' pages 121–140.
% Pass: Opus 5.5 (claude-opus-5-5), 2026-10-03 — first pass, unchecked against the pages by a human.
%
% Page 126 is a cover sheet in his hand (no \page; its words head the section
% that follows). Pages 129 and 130 are bound in reverse order: 130 ends on the
% sentence that 129 opens with.
\documentclass[11pt,a4paper]{article}
\input{../preamble/grothendieck}

\folder{77}
\batch{7}
\pages{121}{140}
\foldertitle{Quadriques affines, extensions quadratiques : notes manuscrites (s.d.).}
\dating{[à partir de 1977]}
\watermark{Édition de démonstration}

\begin{document}

\page{121}
\note{la page s'ouvre au milieu d'une phrase commencée avant ce lot ; la numérotation (13)–(18) continue celle des pages précédentes}
lisse s'exprime par la condition que
\[
  \text{(13)} \qquad \delta'(q) \in \Gamma\bigl(S, (\det T)^{\otimes -2}\bigr)
  \qquad (\text{discr. \struck{divisé}})
\]
\note{le mot « divisé » est souligné puis biffé d'un trait}
soit \uncertain{inversible}, puis par la condition que
\[
  D \xrightarrow{\ \pi \mid D\ } \underline{O}_S
\]
soit un iso, i.e.
\[
  \text{(14)} \qquad T \simeq V \oplus D .
\]
Ceci donne encore une section $a$ canonique de $D$, définie par la condition que
\[
  \text{(15)} \qquad \pi(a) = 1 \qquad (\ill{} \text{ pour } 2\,!)
\]
et on peut donc normaliser $q$ par la condition que
\[
  \text{(16)} \qquad q(a) = 1
\]
de sorte que l'on a encore
\[
  \text{(17)} \qquad q(x,t) = q_0(x) + t^2
\]
et $Q$ se trouve défini par l'équation
\[
  \text{(18)} \qquad q_0(x) + 1 = 0 \quad \text{i.e.} \quad q_0(x) = -1
  \qquad \bigl(x \in \Gamma(-,E)\bigr)
\]
où $E$ est identifié à $V$ grâce à (14).

D'où ici l'équivalence de cat. entre quadriques affines \struck{et} régulières
\add{de rg $n$} \struck{et} formes quadratiques \struck{lisses} de rang $n-1$
\uncertain{munies},

\page{122}
sans restriction à la car. $\neq 2$.

\textbf{2} On va définir une réflexion canonique dans $P$, ayant $H$ comme
hyperplan de pts fixes, et invariant $Q$. Une \add{\uncertain{pareille}}
réflexion \struck{\ill{} d'une autre} ayant $H$ comme hyperplan
\struck{fixe \ill{}} \add{\uncertain{induisant} $E$, donc est associée :
\ill{} automorphisme} \struck{\ill{}} \ill{} de la structure d'extension
\struck{\ill{}} $T$ de $\underline{O}_S$ par $V$, \struck{et} induisant
l'identité sur $\underline{O}_S$, et une homothétie \add{$\lambda\,\mathrm{id}_V$}
sur $V$, \uncertain{ou encore} (quitte à diviser par $\lambda$) par un auto.\ de
la structure d'ext.\ induisant l'identité sur $V$, donc de la forme
\[
  \sigma = \mathrm{id} + v
\]
où $v$ est nulle sur $V$, donc de la forme
\[
  v(x) = \pi(x)\cdot c \qquad \text{i.e.} \qquad v = \pi \otimes c
\]
où $\pi \in \Gamma(S, \check{T})$ est la forme $T \to \underline{O}_S$
définissant la structure d'extension, $c$ une section quelc.\ de $T$.
Écrivons que \struck{\ill{}} $\sigma$ \uncertain{laisse} invariant $q$ mod
homothétie, on trouve
\[
\begin{aligned}
  q \circ \sigma(x) &= q(x + vx) = q\bigl(x + \pi(x)\cdot c\bigr) \\
  &= q(x) + \pi(x)^2 q(c) + \pi(x)\, \varphi(x,c)
\end{aligned}
\]
On veut trouver $\lambda \in \Gamma(S, \mathbb{G}_m)$ tel que
\struck{$q\,u(x) = \lambda q(x)$}
\[
  q\sigma(x) = \lambda q(x) \qquad \forall x \in \Gamma(-,T)
\]
\note{la formule biffée « $q\,u(x) = \lambda q(x)$ » précède, sur la même ligne, la formule retenue}

\page{123}
i.e.
\[
  (\lambda - 1) q = \pi\bigl(q(c)\pi + \psi(c)\bigr)
  \qquad \bigl(\psi(c) \in \Gamma(S, \check{T}),\
  \psi(c)(x) = \varphi(c,x)\bigr)
\]
Comme $q$ est \struck{non dégénérée} \add{lisse} en ttes fibres, ceci montre,
si $\operatorname{rg} T \geqslant 3$, i.e.\ $\operatorname{rg} E \geqslant 2$,
que $\lambda - 1$ est nul en ttes fibres (i.e.\ que $\lambda - 1$ \uncertain{dans}
\uncertain{chaque} \uncertain{nilpotent}), \struck{\ill{}} et \struck{que} comme $\pi$
est \struck{$\neq 0$} \ill{} en ttes fibres, que \uncertain{dans} \ill{}
$\psi(c) + q(c)\pi = 0$ en ttes fibres, … Regardons les solutions
correspondant à $\lambda = 1$, i.e.
\[
  q \circ \sigma = q
\]
qui équivaut à
\[
  \text{(19)} \qquad q(c)\pi + \psi(c) = 0
\]
\struck{\ill{}} On \uncertain{peut} prendre p.\ ex.\ $c = -a$ (dans le cas
\uncertain{$E$} de rg \struck{impair}, qui nous intéresse surtout)
\[
  \text{(20)} \qquad \sigma(x) = x - \pi(x)\, a \qquad (\operatorname{rg} E \text{ pair})
\]
qui est de plus une réflexion i.e.
\[
  \text{(21)} \qquad \sigma^2 = \mathrm{id}_T
\]
(\uncertain{car} \struck{$\sigma^2 =$} $\pi(a) = 1$ …) En fait, c'est la seule
solution \add{\uncertain{vérifiant} $\neq \mathrm{id}$}, \struck{dans} car de (19)
on tire $\psi(c) = \lambda\pi$ \struck{\ill{}} ($\lambda = -q(c)$) i.e.\
$\psi(c) = \lambda\psi(a) = \psi(\lambda a)$ (comme $\psi$ est iso) $c = \lambda a$,
et il faut $\lambda = -q(c)$ i.e.\ $\lambda = -\lambda^2$ i.e.\
$\lambda(\lambda + 1) = 0$, ce qui \uncertain{conduit} \ill{} $\lambda$
\ill{} nul) donc $\lambda = -1$ ….

On dit que $\sigma$ est la symétrie canonique de la quadrique affine régulière.

\marginal{Prop. Il existe une unique réflexion de $P$ admettant $H$ comme
hyperplan de pts fixes et laissant $Q$ \uncertain{invariante}. Sa restriction
à $T$ est donnée par (20) $\sigma x = x - \pi(x)\, a$ i.e.\
$\sigma = \mathrm{id} - \pi \otimes a$}
\note{note marginale écrite verticalement le long du bord gauche}

\page{124}
\note{petit feuillet découpé ; il donne un exemple de la situation des pages précédentes, les matrices $2 \times 2$ avec le déterminant}
\[
\left\{
\begin{aligned}
  T &= M = \left\{ \begin{pmatrix} a & b \\ c & d \end{pmatrix} \right\} \\
  q &= \det = ad - bc
\end{aligned}
\right.
\qquad (b + b',\ c + c')
\]
\note{devant « $\det$ », un mot court biffé, illisible ; « $(b+b', c+c')$ » est une lecture douteuse}
\[
  a = \underline{1} = \begin{pmatrix} 1 & 0 \\ 0 & 1 \end{pmatrix}
\]
\[
  \varphi(a,b,c,d;\, a',b',c',d') = bc' + cb' - (ad' + da')
\]
\note{devant cette formule, un début « $\varphi(x)$ » est biffé}
\[
\begin{aligned}
  \pi(a,b,c,d) &= \varphi(a,b,c,d;\, 1,0,0,1) = -(a+d) \\
  &= -\operatorname{Tr}\begin{pmatrix} a & b \\ c & d \end{pmatrix}
\end{aligned}
\]
attention au signe !
\[
  V = \gamma M = \left\{ \begin{pmatrix} a & b \\ c & -a \end{pmatrix} \right\}
\]
\note{le symbole lu $\gamma$ devant $M$ est douteux ; le sens demande les matrices de trace nulle. Le « $+$ » devant $(a+d)$ est surchargé en « $-$ ». En haut du feuillet, tête-bêche, un fragment coupé par la déchirure : « $\mathbb{Z}/2 \cdot \mathbb{C}P(1) \times \mathbb{C}P(1)$ », lecture douteuse}

\page{125}
\struck{\ill{}} Si $T$ est de rang \add{\uncertain{im}}pair, donc
$T \simeq V \oplus \underline{O}_S$ …, on définit $\sigma$ par
\[
  \text{(22)} \qquad \sigma(x) = x - 2\pi(x)\, a
\]
\note{le numéro (22) est lu sur un chiffre surchargé}
on trouve, sur $E$, par la symétrie par rapport à l'origine. On aura
encore (21).

\section{3. Schéma des automorphismes d'une quadrique affine \add{birégulière}}
\note{titre de sa main, souligné ; « birégulière » est ajouté au-dessus de la ligne}

Soit
\[
  G = \underline{\mathrm{Aut}}(P, H, Q) = \underline{\mathrm{Aut}}_{\mathrm{aff}}(E, Q \cap E)
\]
\struck{$\simeq \mathrm{Aut}(T, \pi, a, q)$}

\struck{$\bigl(\pi \in \Gamma(S, \check{T}),\ a \in \Gamma(S,T),\ q(a) = 1,\ \varphi(a,x) =$}

\note{ces deux lignes sont biffées}
Dans le cas où $T$ de rang impair i.e.\ $E$ de rang pair, on a
\[
  G \simeq \underline{\mathrm{Aut}}(V, q_0) \qquad V \text{ de rang pair}
\]
et désignant par $G^{+}$ le sous-schéma en groupes qui opère trivialement sur
$\det V$, on trouve
\struck{$G \simeq G^{+} \times \mathbb{Z}/2\mathbb{Z}$}
\[
  1 \to G^{+} \to G \to \mathbb{Z}/2\mathbb{Z} \to 1 .
\]
\marginal{$G \simeq \mathrm{Aut}(T, q, a)$}
Si $T$ de rang pair, \struck{\ill{}} on va montrer \uncertain{que} \ill{} un
iso \uncertain{canonique}
\[
  G \simeq G^{0} \times (\mathbb{Z}/2\mathbb{Z})_S
\]
où
\[
  G^{0} \xrightarrow{\ \sim\ } \underline{\mathrm{Aut}}^{+}(V, q_0),
  \qquad (\mathbb{Z}/2\mathbb{Z})_S \text{ engendré par } \sigma .
\]

\section{Formes quadratiques. Fibrés en drapeaux quad., \uncertain{groupes} orthogonaux}
\note{titre pris sur une feuille de garde de sa main (archivistes, p. 126), qui porte aussi « (31 p) » ; la feuille n'a pas d'autre contenu}

\page{127}
\textbf{Formes quadratiques, quadriques et groupes orthogonaux}

\textbf{1. Discriminant et discriminant divisé — formes quadratiques régulières
(ou lisses)}

Si $q$ est une forme quadratique sur $E$ \add{de rang $n$}, on a disc.\ div.
\[
  \text{(1.1)} \qquad \delta'(q) \in \Gamma(\det \check{E})^{\otimes 2}
  = \Gamma(\det E)^{\otimes(-2)}
\]
\marginal{(qui est le disc.\ ordinaire si $n$ pair)}
Conditions équivalentes sur $q$.
\marginal{NB $\delta'(\lambda q) = \lambda^n \delta'(q)$}
\begin{enumerate}
  \item[a)] $\delta'(q)$ inversible
  \item[b)] $V(q) \subset \check{\mathbb{P}}(E)$ est lisse, de codim $1$ en chaque
  fibre (ou encore : est un diviseur relatif lisse)
  \item[c)] $\forall s \in S$, $q \otimes \bar{k}(s)$ \struck{\ill{}} appartient à
  l'unique orbite ouverte de $\underline{\mathrm{Gl}}(E(s))$ opérant sur
  \add{le schéma} \struck{\ill{}} $\underline{\mathrm{Quad}}(E(s))$ (qui est aussi
  l'unique orbite de dimension maximale, laquelle dimension est
  $\frac{n(n+1)}{2}$)
  \item[d)] $\forall s \in S$, $\underline{\mathrm{Aut}}(q \otimes \bar{k}(s))$
  est de dimension minimale (savoir $\frac{n(n-1)}{2}$).
  \item[e')] \add{$\forall s \in S$,} $\exists$ ouvert \add{non vide}
  \struck{\ill{}} $U \subset \underline{\mathrm{Quad}}(E(s))$ tel que
  $q \otimes \bar{k}(s)$ équivalente (sur ext.\ conv.\ de $k(s)$) aux formes
  $\in U$ [NB on peut prendre $U$ défini par $\delta'$ comme fonction pol.\ sur
  $\underline{\mathrm{Quad}}(E)$ …]
  \item[e)] $q$ est \uncertain{isom.} loc ! \struck{\ill{}} [\struck{ou} aussi :
  loc.\ (ét.) si $2$ inv.\ sur $S$ ou $n$ pair] à la forme standard $q_n$
\end{enumerate}
\note{dans d), le « $-1$ » de $\frac{n(n-1)}{2}$ est écrit sur un signe surchargé}

NB Si $q$ est à valeurs dans Mod.\ inv.\ $L$, alors forme bil.\ $\varphi$ est à
val.\ dans $L$, d'où $E \to \check{E} \otimes L$, donc
$\det E \to \det(\check{E} \otimes L) \simeq (\det \check{E}) \otimes L^{\otimes n}$, d'où
\[
  \text{(1.2)} \qquad \delta'(q) \in \Gamma\bigl((\det E)^{\otimes(-2)} \otimes L^{\otimes n}\bigr)
  \simeq \underline{\mathrm{Hom}}\bigl((\det E)^{\otimes 2}, L^{\otimes n}\bigr)
\]

\page{128}
donc si $q'$ est lisse, $\delta'(q)$ définit un iso
\[
  \text{(1.3)} \qquad L^{\otimes n} \simeq (\det E)^{\otimes 2} .
\]

\uncertain{Pour} la suite on suppose \ill{} \add{\uncertain{que} $q$}
\uncertain{spécial-orthogonale} \add{\ill{}}

\textbf{2} Groupe orthogonal\add{, cas $n$ impair $= 2\nu + 1$} // On a
\[
  \text{On pose} \quad
  \left\{
  \begin{aligned}
    &q_n(x_1, y_1, \ldots, x_\nu, y_\nu, z) = \textstyle\sum x_i y_i + z^2 \\
    &\delta'(q_n) = (-1)^\nu
  \end{aligned}
  \right.
\]
\[
  \delta'(-q_n) = -\delta'(q_n) = (-1)^{\nu+1}
\]
\[
\begin{aligned}
  O(q) &= \underline{\mathrm{Aut}}(q) \\
  SO(q) &= \underline{\mathrm{Aut}}(q) \cap Sl(E) = \operatorname{Ker}\bigl(O(q) \xrightarrow{\det} \mathbb{G}_m\bigr)
\end{aligned}
\]
\note{sous $\mathbb{G}_m$, une flèche montante part de $\mu_2$, et une flèche courbe mène de $O(q)$ à $\mu_2$ : le déterminant se factorise par $\mu_2$}

NB Sans condition sur parité de $n$, …
\[
  \text{(2.1)} \qquad \mathbb{G}_m \cap O(q) = \mu_2 \subset O(q)
  \qquad \mu_2 = \underline{\mathrm{Centre}}\, O(q)
\]
\marginal{remarque}
Dans le cas actuel $n$ impair, on a
\[
  \text{(2.2)} \qquad \mu_2 \cap SO(q) = \{1\}
\]

\textbf{Thm} $O(q)$ est le produit direct des ss-groupes $\mu_2$ et $SO(q)$,
$SO(q)$ est lisse sur $S$ à fibres conn.\ (et en fait \struck{\ill{}}
réductifs \add{(i.e.\ $\nu \geqslant 1$)}, \struck{\ill{}} \add{(i.e.\ $\nu \geqslant 3$)}
et semi-simple si $n \geqslant 3$, simple si $n \geqslant 7$).
\note{les deux ajouts entre parenthèses sont écrits sous les mots biffés ; le second est de lecture douteuse}

Structure spéciale orthogonale sur $E$ \add{d'indice $\varepsilon$},
\struck{\ill{}} $\varepsilon$ est une section de $\underline{O}_S^{*}$ (en pratique,
$\varepsilon = \pm 1$) : c'est une section \add{base} $\omega$ de $\det E$, telle que
\[
  \text{(2.3)} \qquad \boxed{\omega^{\otimes -2} = \varepsilon\, \delta'(q)}
\]
On a évidemment
\[
  \text{(2.4)} \qquad \underline{\mathrm{Aut}}(E, q, \omega) = SO(q)
\]
\marginal{Pour $n$ impair, \struck{\ill{}} $(q,\omega) \mapsto (-q,\omega)$
\struck{\ill{}} \uncertain{définit} bijection de l'ens.\ des structures spéciales
orth.\ pour $\varepsilon$, avec celles pour $-\varepsilon$ (dans \ill{}
\uncertain{groupe} \uncertain{adj} …)}
\[
\left\{
\begin{aligned}
  &\text{Si } \varepsilon = (-1)^\nu, \text{ alors } q_n \text{ est un prototype} \\
  &\text{Si } \varepsilon = (-1)^{\nu+1}, \text{ — } -q_n \text{ —}
\end{aligned}
\right.
\]

\page{129}
\note{la page commence au milieu d'une phrase dont le début est au bas de la p. 130 (« l'hom $O'(q) \to \mathbb{G}_m$ induit … ») : les feuillets 129 et 130 se lisent dans l'ordre 130, 129}
$\lambda \mapsto \lambda^2$, sur le centre $\mathbb{G}_m$). Ici pour $n$ impair on
trouve encore

\textbf{Prop} $O'(q) \simeq SO(q) \times \mathbb{G}_m$
\note{sous $SO(q)$ : « $O'(q) \cap Sl(E)$ », précédé d'un mot biffé}
\marginal{\uncertain{commencer} \ill{} \uncertain{par} \ill{} $O(q)$}
qui équivaut au —

\textbf{Cor} (Pour $\varepsilon = \pm 1$ fixé) La catégorie des modules
quasi-orthogonaux \add{$(E, L, q)$ de rang impair} \add{sur $S$} est
équivalente à la cat.\ des \struck{\ill{}} syst.\ $(E_0, q_0, \omega_0, \Delta)$,
où $(E_0, q_0, \omega_0)$ est un module \add{spéc.} orth.\ pour $\varepsilon$, de
rang $n$, $\Delta$ est un module \uncertain{inv.} À un système on associe
\[
  L = \Delta^{\otimes 2}, \quad E = E_0 \otimes \Delta,
  \quad (q_0 \otimes \Delta^{\otimes 2}) : E_0 \otimes \Delta \to \Delta^{\otimes 2} = L,
\]
\[
  \bigl[\text{donc } \det E \simeq \det E_0 \otimes \Delta^{n}
  \simeq \Delta^{2\nu} \otimes \Delta \simeq L^{\nu} \otimes \Delta\bigr]
\]
\note{dans le crochet, les deux dernières expressions sont surchargées et de lecture douteuse}
\marginal{\uncertain{Dire} \uncertain{que} \ill{} plus \ill{}}

Le foncteur quasi-inverse : $(E, L, q)$
\marginal{$n = 2\nu + 1$}
\[
  \Delta = L^{\otimes -\nu} \otimes \det E, \quad E_0 = E \otimes \Delta^{-1},
  \quad q_0 = q \otimes \Delta^{-1} : E_0 = E \otimes \Delta^{-1} \to L\Delta^{-2}
\]
\note{« $L^{\otimes -\nu}$ » est écrit au-dessus de « $\det E$ », sur un mot biffé}
or
\[
  L\Delta^{-2} \simeq L \otimes L^{2\nu} \otimes (\det E)^{-2}
  \simeq L^{n} \otimes (\det E)^{-2} \simeq \underline{O}_S \text{ par (1.3)}
\]
$\omega_0 \in \Gamma(\det E_0)$ déduit de
\[
\begin{aligned}
  \det E_0 &\simeq \det E \otimes \Delta^{-n} \\
  &\simeq \det E\, L^{n\nu} (\det E)^{-n} = L^{n\nu} (\det E)^{-2\nu} \\
  &= \bigl(L (\det E)^{-2}\bigr)^{\nu} \simeq \underline{O}_S^{\otimes \nu} \simeq \underline{O}_S
  \qquad \text{etc.\ …}
\end{aligned}
\]
\note{plusieurs essais biffés entre la première et la deuxième ligne (« $= \Delta\, L^{\nu} \Delta^{-n}$ », « $\det E_0 \simeq \det E \otimes \Delta \ldots$ »). Dans la dernière ligne, la page écrit $L$ et non $L^{n}$ dans la parenthèse}

\page{130}
\note{cette page précède logiquement la p. 129, qui en prolonge la dernière phrase}
\textbf{Prop} (Clas\supplied{sification} des \uncertain{modèles}) Soit
$\varepsilon = \pm 1$. La cat.\ des modules orth.\ \add{$(E,q)$} lisses \add{sur}
$S$ de rg $n = 2\nu + 1$ est équivalente à la cat.\ produit des couples
$(E_0, q_0, \omega_0, \Delta, \varphi)$, où $(E_0, q_0, \omega_0)$ est un module
spécial quadratique pour $\varepsilon$, \struck{\ill{}} $\Delta$ un mod.\ inv.,
$\varphi : \Delta^{\otimes 2} \xrightarrow{\sim} \underline{O}_S$ (NB le couple
$(\Delta, \varphi)$ revient à la donnée d'un torseur sous $\mu_2$). Le foncteur
établissant l'équivalence est
\[
  (E_0, \ldots) \longmapsto \bigl(E_0 \otimes \Delta,\ q_0 \otimes L :
  E_0 \otimes \Delta \to \Delta^{\otimes 2} \xrightarrow[\sim]{\varphi^{-1}} \underline{O}_S\bigr)
\]
\note{une flèche courbe sous la formule relie $E_0 \otimes \Delta$ à $\underline{O}_S$ ; l'exposant de $\varphi$ est lu tel quel}
Le foncteur quasi-inverse est
\[
\begin{aligned}
  (E, q) \longmapsto \bigl(&\Delta = \det E,\ \varphi = \delta'(q)^{-1},\
  E_0 = E \otimes \Delta^{-1}, \\
  &q_0 = \varepsilon q \otimes \Delta^{-1} : E_0 \to \Delta^{-2}
  \xrightarrow{\ {}^t\varphi\ } \underline{O}_S, \\
  &\omega_0 \in \Gamma(\det E_0) \text{ déduit de }
  \det E_0 \simeq (\det E) \otimes \Delta^{-n} \simeq \underline{O}_S\bigr)
\end{aligned}
\]
\note{au-dessus de la dernière ligne, un premier essai biffé : « $\omega_0 \in \Gamma(\det E_0) = \Gamma(\det E) \otimes \Delta^{-n} \simeq \underline{O}_S$ » ; sous $\det E$ et $\Delta^{-n}$, des indications surchargées ($\Delta^{-4}$ ?) de lecture incertaine}

NB $\underline{\mathrm{Aut}}\, V(q) \simeq O(q)/\mu_2 \simeq SO(q)$, il est lisse
connexe \add{réductif} \add{lisse}
\note{sous le NB, deux débuts de ligne biffés, illisibles}

\textbf{(3) Structures de similitude} (formes quadratiques à valeurs dans
loc.\ mod.\ scalaire inv., ou encore forme lisse : valeurs dans un $L$
\[
  q : E \longrightarrow L
\]
Soit $O'(q) \subset Gl(E) \times Gl(L) = Gl(E) \times \mathbb{G}_m$ le groupe de
ses automorphismes, \struck{il \ill{}} \struck{$SO(q)$} des couples $(u, \lambda)$
tels que $q(ux) \overset{\text{idt}}{=} \lambda q(x)$. $O'(q) \to Gl(E)$ est
injectif, $O'(q)$ contient $\mathbb{G}_m$ comme centre (l'hom
$O'(q) \to \mathbb{G}_m$ induit
\note{la phrase se poursuit en haut de la p. 129}

\page{131}
\textbf{(4) Cas $n$ pair $= 2\nu$}
\[
  \text{(4.1)} \qquad q_n = q_{2\nu} = \sum_{1}^{\nu} (x_i y_i)
  \qquad \delta'(q_n) = (-1)^\nu
\]
\struck{(4.} Ici on a
\[
  \text{(4.2)} \qquad \mu_2 \subset SO(q) \cap Sl(E)
\]
\note{sous $\mu_2$ : « $= \mathrm{Centre}\, O(q)$ » ; sous $SO(q) \cap Sl(E)$, d'abord écrit avec un $S$ surchargé : « on ne l'appelle [mot biffé] $SO(q)$ à cause des car.\ $2$ »}

\textbf{Thm} Pour $n$ pair, $O(q)$ est lisse, $SO(q)$ est sa composante neutre,
on a
\[
  O(q)/SO(q) \simeq (\mathbb{Z}/2)_S
\]
et
\[
  \mu_2 = \underline{\mathrm{Centre}}\bigl(O(q)\bigr) \subset SO(q) .
\]
\note{au-dessus de « Thm », un mot biffé}

\textbf{Cor} Si $X = V(q)$, le groupe $\underline{\mathrm{Aut}}(X) \simeq O(q)/\mu_2$
est lisse, et
\[
\left\{
\begin{aligned}
  &\underline{\mathrm{Aut}}(X)^{0} \simeq SO(q)/\mu_2 \\
  &\underline{\mathrm{Aut}}(X)/\underline{\mathrm{Aut}}(X)^{0} \simeq (\mathbb{Z}/2)_S
\end{aligned}
\right.
\]

Indications de dém.\ du Thm.

La chose essentielle est de définir l'hom \struck{\ill{}} \add{épi}
\[
  O(q) \longrightarrow (\mathbb{Z}/2)_S
\]
et pour ceci on va associer canon.\ à $q$ un rev.\ étale de degré $2$ de $S$.

\page{132}
\textbf{Prop} Considérons le sous-schéma de $\underline{\mathrm{Grass}}_\nu(E)$
\struck{formé} \add{correspondant} \struck{aux} « sous-espaces isotropes
\struck{maximaux} » [sur lesquels $q$ \struck{et la forme bil.\ associée}
\add{est} nuls ; $q \mid V \equiv 0$,] \struck{$\varphi \mid V \times E \equiv 0$}.
C'est un sous-schéma fermé, il est lisse, et \struck{\ill{}} \uncertain{pour}
les fibres géom., le nb de composantes connexes \struck{est} $2$.
\note{le reste de la page est blanc}

\section{Fibrés en drapeaux associés à \add{une} \struck{des} formes quadratiques régulières}

\page{133}
\note{titre de sa main, en tête de la page}
$(E,q)$, avec $E$ de rg $N = 2n+1$ (cas impair) — \add{\uncertain{ou}} $2n$,
$n \geqslant 1$ \add{donc $N \geqslant 2$}). Soit $H \subset [1,n] \cap \mathbb{Z}$,
[donc $H = \{i_1, \ldots, i_\nu\}$, où $1 \leqslant i_1 < \cdots < i_\nu \leqslant n$]
\struck{On pose} $\nu = \operatorname{card} H$.
\[
  \Delta_H(E,q) = \text{ens.\ des drapeaux de type } H \text{ de } E
\]
\[
  E_1 \subset E_2 \subset \cdots \subset E_\nu
  \qquad \bigl(\nu = \operatorname{card} H,\ H = \{\operatorname{rg} E_1, \ldots, \operatorname{rg} E_\nu\}\bigr)
\]
tels que \struck{les} $E_\nu$ soit tot\supplied{alement} isotrope.
\[
  \underline{\Delta}_H(E,q) : \mathrm{Sch}_{/S} \longrightarrow \mathrm{Ens}
  \qquad S' \longmapsto \Delta_H(E_{S'}, q_{S'})
\]

\textbf{Théorème !} a) $\underline{\Delta}_H(E,q) \hookrightarrow \underline{\mathrm{Drap}}_H(E)$
est \struck{représentable} \add{une} immersion fermée — en particulier
$\underline{\Delta}_H(E,q)$ est représentable, et projectif sur $S$.

b) $\underline{\Delta}_H(E,q)$ est lisse sur $S$. Il a ses fibres géométriquement
$0$-connexes, sauf pour \struck{$i_\nu = n$}, rg $E$ pair $= 2n$, $i_\nu = n$,
auquel cas \struck{\ill{}} ses fibres géom.\ ont exactement deux composantes
connexes.

\textbf{Dém} Le fait que $\underline{\Delta}_H(E,q) \to \underline{\mathrm{Drap}}_H(E)$
est une imm.\ fermée est immédiat, …

Pour la lissité de $\underline{\Delta}_H(E,q)$ \struck{$(F,q)$ est une}, on utilise
le critère infinitésimal de lissité — on est ramené à prouver que si
$S_0 \subset S$ est défini par $J$ idéal \uncertain{quasi-coh.}\ de carré nul,

\marginal{Donner aussi la structure de $\underline{\Delta}_H \to \underline{\Delta}_K$
pour $K \subset H$ … (si $K \neq \emptyset$, \ill{} $\underline{\Delta}_H \to S$ …)
$\underline{\Delta}_H \to \underline{\Delta}_K$ est lisse \uncertain{projectif}, ses
fibres \uncertain{géom.}\ sont $0$-conn., \uncertain{sauf} si \struck{\ill{}}
$i_\nu = \operatorname{Max} H = n$, \uncertain{$\nu$} \ill{} rg $E = 2n$,
\uncertain{auquel} cas les fibres géom.\ ont $2$ comp.\ conn.}
\note{note marginale écrite en oblique le long du bord gauche}

\page{134}
$S$ affine, et si $(E_{1,0}, \ldots, E_{\nu,0})$ est un élément de
$\Delta_H(E_0, q_0)$, celui-ci se remonte \add{en un
$(E_1, \ldots, E_\nu) \in \Delta_H(E,q)$}. Il suffit bien sûr (\struck{par la}
\add{utilisant} propriétés de lissité de
$\underline{\mathrm{Drap}}_H E \to \underline{\mathrm{Grass}}_{i_\nu}(E)$) de prouver
que $F_0 = E_{\nu,0}$ se remonte en un sous-Module $F = E_\nu$ \struck{\ill{}}
f.d.\ direct et tot.\ isotrope.

\note{tout ce qui suit sur la page est annulé par de longs traits obliques ; on en donne ce qui se lit}
\struck{On considère \ill{} \ill{} sur $\operatorname{rg} F$. Si $d = 1$ \ill{}
remonte \ill{} $F$} \struck{on \ill{}} ($\underline{\mathrm{Grass}}_{i_\nu}(E)$ lisse
sur $S$) \struck{facteur direct, mais $q \mid F$ est à valeurs dans $J$, donc peut
s'identifier à un él.\ de $\mathrm{Quad}(F_0, J)$ par nécessité,
\ill{} une « connection » sur le choix de $F$ dépend de}
\[
  \text{\struck{$\Gamma\,\underline{\mathrm{Hom}}(F_0, E_0/F_0 \otimes J)
  \simeq \Gamma\bigl(\underline{\mathrm{Hom}}(F_0, E_0/F_0) \otimes J\bigr)$}}
\]
\struck{puis supposons que $F$ a une base $(e_{i,0})_{1 \leqslant i \leqslant d}$
($d = i_\nu$). On \add{note} que chaque $e_{i0}$ peut se remonter en $e_i$
\add{$e \in E$ tel que $q(e_i) = 0$} (lissité de la quadrique projective définie
par $q$), sans}
\struck{$F$ le module engendré par les} \struck{c'est la lissité de la quadrique
projective. Si $d > 1$, \ill{} $F_0 = G_0 + \underline{O}_S e_0$, $G_0$ de rang
$d-1$, $e$ section de $F_0$. On remonte $G_0$ en $G \in \Delta_{\{d-1\}}(E,q)$,
puis il reste à remonter $e_0$ en une section $e$ de $E$}
\note{cette phrase annulée se poursuit en haut de la p. 135}

\page{135}
\note{la moitié supérieure de la page est annulée par de grands traits en croix}
\struck{satisfaisant}
\[
  \text{\struck{$q(e) = 0,$}} \quad \text{\struck{$\varphi(e, G) = 0$}}
\]
\struck{Commençons par remonter en $e$ satisfaisant la première de ces relations.
Alors}
\[
  \text{\struck{$x \mapsto \varphi(e, x) \quad G \to \underline{O}_S$}}
\]
\struck{est nul mod $J$ donc définit}
\[
  \text{\struck{$\psi_0 : G_0 \to J$}} \qquad
  \text{\struck{$\psi_0 \in \underline{\mathrm{Hom}}(G_0, J)$}}
\]
\struck{Essayons de corriger $e$ par}
\[
  \text{\struck{$e' = e + \xi$}} \qquad \text{\struck{$\xi \in \Gamma(JE)$}}
\]
\struck{la condition $q(e') = 0$ s'écrit, compte tenu que $q(e) = 0$,}
\[
  \text{\struck{$q(\xi) + \varphi(e, \xi) = 0$}}
\]
\struck{or encore $q(\xi) = 0$}

Considérons l'hom
\[
\begin{aligned}
  \underline{\mathrm{Hom}}(F_0, E_0/F_0) \otimes J &\longrightarrow \underline{\mathrm{Quad}}(F_0) \\
  u_0 &\longmapsto \bigl(x_0 \mapsto \varphi_0(x_0, u_0 x_0)\bigr)
\end{aligned}
\]
\note{un symbole surchargé suit la parenthèse de $\underline{\mathrm{Hom}}(F_0, E_0/F_0)$ ; le $\otimes J$ est une lecture}
[NB $\varphi(F_0, F_0) = 0$, donc $\varphi \mid F_0 \times E_0$ se factorise par
$F_0 \times E_0/F_0$], il suffit de prouver qu'il est \struck{surjectif} — il revient
au même, par dévissage par $J$, de le prouver pour sections sur base affine …
Or c'est un composé
\[
  \underline{\mathrm{Hom}}(F_0, E_0/F_0) \xrightarrow{\ \text{épi}\ }
  \underline{\mathrm{Hom}}(F_0, E_0/F_0^{\perp}) \xrightarrow[\text{épi}]{\ \sim\ }
  \underline{\mathrm{Quad}}(F_0)
\]
\note{au-dessus de $\underline{\mathrm{Hom}}(F_0, E_0/F_0^{\perp})$ : « $\simeq \check{F}_0$, isom : explicitement ($q$ régulière !) cf.\ plus haut » ; dessous, une première version biffée de la même indication}
\marginal{NB}

\page{136}
Il faut rajouter le

\textbf{Lemme 1} Soit $F \subset E$, \struck{\ill{}} \add{$F \in \Delta_{\{d\}}(E,q)$},
alors $F^{\perp}$ est loc.\ f.d.\ direct et $\supset F$, et l'accouplement entre
$F$ et $E/F^{\perp}$ est une dualité (donc $F = F^{\perp\perp}$, la formation de
$F^{\perp}$ commute au changement de base …)
\add{$F^{\perp}$}

Il faut regarder \struck{\ill{}} le composé $E \to \check{E} \to \check{F}$ et
montrer que c'est \struck{surjectif} \add{épi} — le noyau étant $F^{\perp}$, le
résultat s'ensuit. Or on est ramené au cas où $S$ est spectre d'un corps
alg.\ clos [le seul cas non \uncertain{évident} est celui de qu.\ dégénérée, i.e.\
car $k = 2$, rg $E$ impair]. Mais \struck{\ill{}, \ill{} $E \to \check{E}$ est
\ill{}} si $a \in F$ est orth.\ à l'image de $E$, il est orth.\ à $E$ par $\varphi$,
\struck{\ill{}} ce qui joint à $q(a) = 0$ implique $a = 0$, cqfd.

Reste : vérifier les assertions de connexion. Par dévissage on est ramené, pour le
cas \struck{\ill{}} rg $E$ impair, ou rg $E$ pair, $i_\nu < n$,

\textbf{Lemme 2} Soient $0 \leqslant d < d' \leqslant n$,
\struck{\ill{}} \struck{$d' < n$ si rg $E$ pair}. Soit $F \in \Delta_{\{d\}}(E,q)$, considérons
\note{au-dessus de « $0 \leqslant d$ » : « NB on admet ici $d = 0$ » ; avant « Lemme 2 », un mot biffé}

\page{137}
\uncertain{divers}, le sous-schéma de $\underline{\Delta}_{\{d,d'\}}(E,q)$
\struck{défini par les} \struck{$(F,G)$} \add{image inverse de la section $F$ par}
\[
  \underline{\Delta}_{\{d,d'\}}(E,q) \longrightarrow \underline{\Delta}_{d}(E,q) .
\]
Le schéma est \add{(lisse projectif),} \struck{à fibres non vides connexes}
\add{à fibres géom.\ $0$-connexes} sauf si rg $E$ pair et $d' = n$, auquel cas elles
ont $2$ comp.\ conn.
\struck{Dém.\ \ill{} \ill{}, \ill{} $d' = d + 1$ (\ill{} (de \ill{} \ill{} \ill{} …))}
\note{le lemme 2 est encadré d'un trait ; une première démonstration, ébauchée sous l'énoncé, est biffée et entourée}

Considérons $F^{\perp}$, (cf.\ Lemme 1) et considérons,
\struck{\ill{} $S$ supposé affine, un scindage}
\[
  \text{\struck{$F^{\perp} \simeq F \oplus G'$}} \qquad E' = F^{\perp}/F
\]
\add{muni de $q'$ déduit de $q$ (bien \uncertain{posé})}.
\struck{soit} \textbf{Lemme 3} \struck{\ill{}} \struck{$E \simeq G \oplus G^{\perp}$}
\struck{La forme} $E'$ est loc.\ libre \struck{et $G^{\perp}$ loc.} de rg $N - 2d$
(donc on a \ill{} vérifié que le rg de $E$) et $q'$ est régulière. Les
$G \in \Delta_{\{d'\}}(E,q)$ tels que $F \subset G$ correspondent donc
biunivoquement aux $G' \in \Delta_{\{d'-d\}}(E')$.

Ceci nous ramène à déterminer les composantes connexes d'un $\Delta_{\{d\}}(E)$.
Ici encore, par le cas de la \uncertain{connexité}, on est ramené par dévissage au
cas de $d = 1$, donc \add{\uncertain{de} rg $E$} \add{: fibres géom.}

\textbf{Lemme 4} $\Delta_{\{1\}}(E)$ est connexe si rg $E \geqslant 3$, ses fibres
géom.\ ont $2$ \uncertain{cc.}\ si rg $E = 2$.
Évident \struck{car} \uncertain{une} \uncertain{conique} dans $\mathbb{P}(E)$ …
\note{l'ajout « : fibres géom. » au-dessus de la ligne précédente semble destiné à l'énoncé du lemme 4 (« à fibres géom.\ connexes »)}

\page{138}
\struck{(\ill{})}

Reste le cas $N = 2n$, \struck{$H$} $H = \{n\}$ i.e.\ $\Delta_{\{n\}}(E,q)$.
\struck{Soit} Supposons d'abord $2$ inv.\ sur $S$. Dans $F \in \Delta_{\{n\}}(E,q)$,
\struck{\ill{}} donc $F = F^{\perp}$ (via Lemme 1), \struck{\ill{}}
$E/F \xrightarrow{\sim} \check{F}$, d'où un isom.\ can.
\[
  \underbrace{\det F \otimes \det \check{F}}_{\simeq\, \underline{O}_S} \longrightarrow \det E
\]
\note{devant la formule, un mot biffé}
qu'on interprète comme une section de $\det E$, notée $\omega_F$. Choisissons un
supplémentaire \add{$G$} (S supposé affine) de $F$ dans $E$, donc
\[
  E \simeq F \oplus G \qquad \varphi \mid F \times G \text{ une dualité}
\]
Soit $(e_1, \ldots, e_n)$ une base de \uncertain{$E$} (\ill{} loc.\ (Zar)),
$e_1^{*}, \ldots, e_n^{*}$ la base duale \struck{\ill{}} de $G \simeq \check{F}$,
définie par
\[
  \varphi(e_i, e_j^{*}) = \delta_{ij}
\]
on a donc
\[
  \omega_F = e_1 \wedge \cdots \wedge e_n \wedge e_1^{*} \wedge \cdots \wedge e_n^{*}
\]
\note{« base de $E$ » se lit ainsi ; le sens demande une base de $F$}

\page{139}
[\textbf{Lemme 5} (\uncertain{cas} \uncertain{général}) Soit $F \in \Delta_{\{n\}}(E,q)$, où
$N = \operatorname{rg} E$ \struck{\ill{}} $= 2n$. Alors $\exists$ suppl.\ $G$ de $F$ tel
que $q \mid G = 0$. (La forme $q$ sur $E \simeq F \oplus G$ s'exprime alors par
\[
  q(x \oplus y) = \varphi(x,y) \qquad
  \bigl(= \langle x, y\rangle \text{ si on identifie } G \text{ à } \check{F}\bigr)
\]

\textbf{Dém} Car choisit un suppl.\ quelc., puis on le corrige, en utilisant que
\[
  \underline{\mathrm{Hom}}(E/F, F) \longrightarrow \underline{\mathrm{Quad}}(G)
  \text{ est surjectif …}]
\]
\note{sous $E/F$ et $F$, des indications d'identification ($\simeq \check{F}$ ? et $\simeq \check{G}$), lecture douteuse}

Donc dans le raisonnement précédent, OPS $q \mid G = 0$, de sorte que le calcul du
discriminant donne
\[
  \delta(q) = (-1)^n (\omega_F)^2
\]
on en tire
\[
  \boxed{(\omega_F)^2 = (-1)^n \delta(q)^{-1}}
\]
\note{les deux formules donnent $\delta(q)$ puis $\delta(q)^{-1}$ ; on transcrit tel quel}

\struck{Considérons d'autre part} On définit par $F \mapsto \omega_F$ un morphisme
\add{dit \uncertain{des} \uncertain{vecteurs} d'Arf …}
\[
  \underline{\Delta}_{\{n\}}(E,q) \longrightarrow M\bigl((-1)^n \delta(q)\bigr)
  = \underline{\mathrm{Arf}}(E,q)
\]
[$=$ sous-schéma de $V(\underline{\det}\, E)$ des racines carrées de
$(-1)^n \delta(q)^{-1}$]

\page{140}
Nous prouvons maintenant ici \add{et supposons $2$ inv.\ sur $S$.}

\textbf{Thm 2} Supposons rg $E = N$ pair $= 2n$. Alors le morphisme précédent est
lisse projectif à fibres \add{géom.} $0$-connexes.

Lisse projectif est évident, il faut voir que les fibres \add{géom.} sont
$0$-connexes. \add{(sur $k$ alg.\ clos)}
\struck{Elles sont non vides : il suffit (\ill{} \ill{} \ill{} $F'$ \ill{}
$\omega_{F'} = -\omega_F$. Il suffit pour ceci de prendre prendre}
\[
\begin{aligned}
  &\text{\struck{$F' = \underline{O}_S e_1^{*} + \sum_{2 \leqslant i \leqslant n} \underline{O}_S e_i$}} \\
  &\text{\struck{$G' = \underline{O}_S e_1 + \sum_{2 \leqslant i \leqslant n} \underline{O}_S e_i^{*}$}}
\end{aligned}
\]
\note{ce passage est encadré et barré de traits obliques}
Le fait que \uncertain{ce} \ill{} \uncertain{est} surjectif provient du fait que
$\underline{\mathrm{Aut}}(E,q) = O(q)$ opère transitivement sur
$\underline{\mathrm{Arf}}(E,q)$ (une réflexion orthogonale p.\ ex.\ \struck{\ill{}}
\add{échange les deux points} \struck{\ill{}} du rev.\ d'Arf) il reste à prouver que
les \add{\uncertain{schémas}} \add{$\underline{\Delta}_{\{n\}}(E,q)$} $F$ tels que
$\omega_F = \omega$ (fixé) \struck{forment} forment un schéma connexe (sur le
\uncertain{rg.} clos). Or c'est l'image des $\underline{\Delta}_{\{1, \ldots, n\}}$, où
$H = [1,n]$, donc on est
\note{l'indice de $\underline{\Delta}$ est surchargé ; la phrase se poursuit au-delà de ce lot}

\end{document}
